\documentclass[11pt]{article}
\usepackage{mn}

\begin{document}
\zestaw{1}{Operacje na „prawdziwych” liczbach: arytmetyka zmiennoprzecinkowa, utrata cyfr znaczących, norma macierzy}

\begin{teoria}[Jak czytać ten dokument]
Każde zadanie ma: przepisaną \textbf{treść} (szare pudełko), potrzebną
\textbf{teorię} (niebieskie), \textbf{intuicję} (żółte), \textbf{rozwiązanie
formalne} i~\textbf{odpowiedź} (zielone). Wyniki zostały sprawdzone
numerycznie skryptem \texttt{weryfikacja/mn01.py} (Python liczy na
\texttt{float}, czyli dokładnie w~opisanym niżej formacie \emph{double}).
\end{teoria}

% =====================================================================
\section{Teoria: liczby zmiennoprzecinkowe}

\begin{definicja}[Liczba zmiennoprzecinkowa, IEEE 754 \emph{double}]
Liczba maszynowa (poza zerem, nieskończonościami i~liczbami
zdenormalizowanymi) ma postać
\[
x=\pm(1.d_1d_2\ldots d_{52})_2\cdot 2^{e},\qquad e\in\{-1022,\ldots,1023\}.
\]
Część $1.d_1\ldots d_{52}$ to \emph{mantysa} (53 bity znaczące, pierwszy
zawsze $1$), $e$ to \emph{cecha} (wykładnik). Największa liczba to
$(2-2^{-52})\cdot2^{1023}\approx1{,}8\cdot10^{308}$.
\end{definicja}

Ważne konsekwencje:
\begin{itemize}
\item W~przedziale $[2^e,2^{e+1})$ liczby maszynowe są rozmieszczone
\textbf{równomiernie} co $2^{e-52}$ (tzw.\ $\mathrm{ulp}$, \emph{unit in the
last place}). Np.\ w~$[1,2)$ co $2^{-52}$, w~$[2,4)$ co $2^{-51}$,
w~$[2^{53},2^{54})$ co $2$ --- liczby $2^{53}+1$ już nie da się zapisać.
\item Każda liczba całkowita o~module $\le2^{53}$ oraz każda liczba postaci
$m\cdot2^{k}$ z~$\abs{m}<2^{53}$ (w~zakresie wykładników) jest maszynowa.
Liczba $0{,}1$ \textbf{nie} jest maszynowa (ma nieskończone rozwinięcie
dwójkowe $0{,}0\overline{0011}_2$).
\end{itemize}

\begin{definicja}[Zaokrąglenie i działania maszynowe]
$\fl(x)$ to liczba maszynowa najbliższa $x$; przy remisie (gdy $x$ leży
dokładnie w~połowie między dwiema sąsiednimi) wybieramy tę z~\textbf{parzystą}
ostatnią cyfrą mantysy (\emph{round half to even}). Działania maszynowe to
\[
x\oplusm y=\fl(x+y),\qquad x\otimesm y=\fl(x\cdot y),
\]
czyli: policz wynik dokładnie, potem zaokrąglij --- tak działa IEEE 754.
\end{definicja}

\begin{twierdzenie}[Standardowy model błędu]\label{tw:model}
Jeśli wynik nie wychodzi poza zakres (brak nadmiaru i~niedomiaru), to
\[
\fl(x\circ y)=(x\circ y)(1+\delta),\qquad \abs{\delta}\le\eps=2^{-53}
\approx1{,}1\cdot10^{-16}
\]
dla $\circ\in\{+,-,\cdot,/\}$, a~także $\fl(\sqrt x)=\sqrt x\,(1+\delta)$.
\end{twierdzenie}

\begin{proof}
Niech $z=x\circ y\in[2^e,2^{e+1})$ (dla ujemnych symetrycznie). Sąsiednie
liczby maszynowe są tu odległe o~$2^{e-52}$, więc najbliższa jest w~odległości
co najwyżej połowy: $\abs{\fl(z)-z}\le2^{e-53}$. Stąd
$\abs{\delta}=\abs{\fl(z)-z}/\abs z\le 2^{e-53}/2^{e}=2^{-53}$.
\end{proof}

\begin{uwaga}
Wielkość $\eps=2^{-53}$ to \emph{precyzja arytmetyki} (\emph{unit roundoff});
często nazywa się też \emph{epsilonem maszynowym} liczbę $2^{-52}$ --- odległość
od $1$ do następnej liczby maszynowej. Różne książki używają różnych konwencji.
\end{uwaga}

\begin{intuicja}
Każde pojedyncze działanie jest „prawie dokładne” --- względny błąd
$\le10^{-16}$. Ale działania maszynowe to \emph{inne działania} niż
zwykłe $+$ i~$\cdot$: po każdym kroku coś obcinamy, więc kolejność kroków
ma znaczenie. Prawa algebry (łączność, rozdzielność) przestają obowiązywać,
a przemienność zostaje (bo $x+y$ i~$y+x$ to ta sama liczba przed
zaokrągleniem).

Najprościej zobaczyć to w~„zabawkowym” systemie dziesiętnym z~3~cyframi
znaczącymi: $1{,}00+0{,}004=1{,}004\to1{,}00$, więc
$(1{,}00+0{,}004)+0{,}004=1{,}00$, a~$1{,}00+(0{,}004+0{,}004)=1{,}008\to1{,}01$.
Małe składniki dodawane osobno „giną”, a~dodane najpierw do siebie ---
przeżywają. W~\emph{double} robimy dokładnie to samo, tylko w~systemie
dwójkowym z~53 cyframi.
\end{intuicja}

% =====================================================================
\section{Zadanie 1. Dodawanie nie jest łączne}
\begin{tresc}
Pokaż, że dodawanie liczb zmiennoprzecinkowych nie jest łączne --- znajdź
takie liczby $a,b,c$, dla których przy dodawaniu maszynowym
\[
(a+b)+c\ne a+(b+c).
\]
Same liczby $a,b,c$ muszą dać się dokładnie zapisać, a~dodatkowo wszystkie
liczby oraz ich suma mają zmieścić się w~zakresie największej możliwej do
zapisania liczby.
\end{tresc}

\begin{intuicja}
Warunki z~treści wykluczają dwa „tanie” przykłady:
\begin{itemize}
\item $0{,}1+0{,}2+0{,}3$ --- tu różnica bierze się częściowo stąd, że
$0{,}1$, $0{,}2$, $0{,}3$ w~ogóle nie są liczbami maszynowymi;
\item $(M+M)-M=\infty\ne M=M+(M-M)$ dla największej liczby $M$ --- to
nadmiar, a~nie błąd zaokrąglenia.
\end{itemize}
Chcemy więc zobaczyć \emph{czysty} efekt zaokrąglenia. Robimy to jak
w~systemie zabawkowym: do $1$ dodajemy liczbę równą \emph{połowie} odstępu
między $1$ a~następną liczbą maszynową. Raz --- remis, zaokrąglenie do
parzystej, wynik wraca do $1$. Dwa razy osobno --- dwa razy wraca do $1$.
Ale dwie połówki dodane najpierw do siebie dają cały odstęp, który już
„widać”.
\end{intuicja}

\formalnie
Niech $a=1$, $b=c=2^{-53}$. Wszystkie trzy są liczbami maszynowymi
(potęgi dwójki z~zakresu), sumy są rzędu $1$, więc daleko od nadmiaru.

\emph{Lewa strona.} Najbliższe liczbie $1+2^{-53}$ liczby maszynowe to $1$
i~$1+2^{-52}$ (odstęp w~$[1,2)$ wynosi $2^{-52}$), a~$1+2^{-53}$ leży
dokładnie pośrodku. Remis rozstrzygamy na korzyść parzystej mantysy:
$1=(1.00\ldots0)_2$ ma ostatni bit $0$, a~$1+2^{-52}=(1.00\ldots01)_2$ --- bit
$1$. Zatem
\[
a\oplusm b=\fl(1+2^{-53})=1,\qquad (a\oplusm b)\oplusm c=\fl(1+2^{-53})=1.
\]

\emph{Prawa strona.} $b\oplusm c=\fl(2^{-52})=2^{-52}$ (dokładnie), więc
\[
a\oplusm(b\oplusm c)=\fl(1+2^{-52})=1+2^{-52},
\]
bo $1+2^{-52}$ jest liczbą maszynową. Wyniki się różnią:
\[
(a\oplusm b)\oplusm c=1\ne1+2^{-52}=a\oplusm(b\oplusm c).
\]

\begin{uwaga}
Ten sam mechanizm w~liczbach całkowitych: $a=2^{53}$, $b=c=1$.
W~$[2^{53},2^{54})$ odstęp to $2$, więc $2^{53}+1$ to remis
i~$\fl(2^{53}+1)=2^{53}$, stąd $(a\oplusm b)\oplusm c=2^{53}$, ale
$a\oplusm(b\oplusm c)=2^{53}+2$. (To jest powód, dla którego np.\ w~JavaScripcie
\texttt{2**53 + 1 === 2**53}.)
\end{uwaga}

\begin{odpowiedz}
W~\emph{double}: $a=1$, $b=c=2^{-53}$. Wtedy
$(a+b)+c=1$, a~$a+(b+c)=1+2^{-52}$ --- dodawanie maszynowe nie jest łączne.
\end{odpowiedz}

% =====================================================================
\section{Zadanie 2. Mnożenie nie jest rozdzielne względem dodawania}
\begin{tresc}
To samo co zadanie~1, ale dla rozdzielności mnożenia względem dodawania:
znajdź $a,b,c$ (dokładnie zapisywalne, bez wychodzenia poza zakres), dla
których przy działaniach maszynowych
\[
a\cdot(b+c)\ne a\cdot b+a\cdot c.
\]
\end{tresc}

\begin{intuicja}
Ten sam trik: $b+c$ ma „zgubić” mały składnik $c$, ale po pomnożeniu przez
$a$ ten składnik ma być już na tyle duży, żeby przeżył dodawanie $ab+ac$.
Mnożenie przez $a=3$ zmienia skalę: $3$ leży w~$[2,4)$, gdzie odstęp to
$2^{-51}$, a $3c=\tfrac34\cdot2^{-51}$ jest \emph{większe niż pół odstępu},
więc zaokrągla się w~górę. Uwaga: $a$ nie może być potęgą dwójki ---
mnożenie przez $2^k$ jest dokładne i~niczego nie zmienia.
\end{intuicja}

\formalnie
Niech $a=3$, $b=1$, $c=2^{-53}$ (wszystkie maszynowe, wyniki rzędu $3$).

\emph{Lewa strona.} Jak w~zadaniu~1: $b\oplusm c=\fl(1+2^{-53})=1$ (remis,
do parzystej), więc $a\otimesm(b\oplusm c)=\fl(3\cdot1)=3$.

\emph{Prawa strona.} $a\otimesm b=3$ oraz $a\otimesm c=\fl(3\cdot2^{-53})=3\cdot2^{-53}$
(dokładnie: $3\cdot2^{-53}$ ma mantysę $(1.1)_2$). Teraz
\[
3+3\cdot2^{-53}=3+\tfrac34\cdot2^{-51}.
\]
W~$[2,4)$ sąsiednie liczby maszynowe to $3$ i~$3+2^{-51}$. Odległość od $3$
to $\tfrac34\cdot2^{-51}$, od $3+2^{-51}$ --- $\tfrac14\cdot2^{-51}$, więc
\[
(a\otimesm b)\oplusm(a\otimesm c)=3+2^{-51}\ne3=a\otimesm(b\oplusm c).
\]

\begin{uwaga}
Warunek „bez wychodzenia poza zakres” wyklucza np.\ $a=2$, $b=M$, $c=-M$
($M$ --- największa liczba): $2(M-M)=0$, ale $2M=\infty$ i~$\infty-\infty=\mathrm{NaN}$.
\end{uwaga}

\begin{odpowiedz}
$a=3$, $b=1$, $c=2^{-53}$: $\;a(b+c)=3$, a~$ab+ac=3+2^{-51}$.
\end{odpowiedz}

% =====================================================================
\section{Zadanie 3. Pierwiastki równania kwadratowego bez \emph{catastrophic cancellation}}
\begin{tresc}
Dane są $A,B,C\in[0,20]$, gdzie $B>10$. Zaproponuj algorytm obliczania
pierwiastków równania kwadratowego $Ax^2+Bx+C=0$, w~którym nie wystąpi
zjawisko \emph{catastrophic cancellation}.
\end{tresc}

\begin{teoria}[Utrata cyfr znaczących]
\begin{definicja}[Błąd względny]
Jeśli $\tilde x$ przybliża $x\ne0$, to błąd względny to
$\abs{\tilde x-x}/\abs x$. Błąd względny $10^{-k}$ $\approx$ „$k$ poprawnych
cyfr znaczących”.
\end{definicja}

\begin{lemat}[Odejmowanie bliskich liczb]\label{lem:cancel}
Niech $\tilde x=x(1+\delta_1)$, $\tilde y=y(1+\delta_2)$, $\abs{\delta_i}\le\delta$.
Wtedy
\[
\frac{\abs{(\tilde x-\tilde y)-(x-y)}}{\abs{x-y}}
=\frac{\abs{x\delta_1-y\delta_2}}{\abs{x-y}}
\le\delta\cdot\frac{\abs x+\abs y}{\abs{x-y}}.
\]
Natomiast dla $x,y$ \textbf{tego samego znaku}
\[
\frac{\abs{(\tilde x+\tilde y)-(x+y)}}{\abs{x+y}}
\le\delta\cdot\frac{\abs x+\abs y}{\abs{x+y}}=\delta.
\]
\end{lemat}

Czynnik $\frac{\abs x+\abs y}{\abs{x-y}}$ jest ogromny, gdy $x\approx y$:
samo odejmowanie liczone jest (prawie) dokładnie, ale wzmacnia błędy, które
\emph{już siedziały} w~$\tilde x,\tilde y$. To jest \textbf{catastrophic
cancellation}: wiodące (poprawne) cyfry się skracają, zostają same
„śmieciowe”. Dodawanie liczb tego samego znaku błędu nie wzmacnia.
\end{teoria}

\begin{intuicja}
Szkolny wzór $x_{1,2}=\frac{-B\pm\sqrt{\Delta}}{2A}$ ma pułapkę: gdy
$4AC\ll B^2$, to $\sqrt\Delta\approx B$ i~w~$-B+\sqrt\Delta$ odejmujemy dwie
prawie równe liczby. Przykład: $A=C=10^{-9}$, $B=20$. Wtedy
$\Delta=400-4\cdot10^{-18}$ zaokrągla się do $400$, $\sqrt\Delta=20$
i~wzór daje $x_1=0$, a~prawdziwy pierwiastek to $\approx-5\cdot10^{-11}$
--- \emph{zero} poprawnych cyfr.

Ratunek: liczymy tylko ten pierwiastek, w~którym nic się nie skraca
(ponieważ $B>0$, jest to ten ze znakiem „$-$”: $-B-\sqrt\Delta$, suma dwóch
liczb ujemnych), a~drugi bierzemy ze wzorów Viète'a $x_1x_2=C/A$ --- dzielenie
też nie skraca cyfr.
\end{intuicja}

\formalnie
\begin{algorithm}[H]
\caption{Pierwiastki $Ax^2+Bx+C=0$, $B>10$}
\begin{algorithmic}[1]
\State $\Delta\gets B\cdot B-4\cdot A\cdot C$
\If{$\Delta\ge0$}
  \State $q\gets-\bigl(B+\sqrt\Delta\bigr)/2$ \Comment{$q\le-5$}
  \State $x_2\gets C/q$
  \If{$A=0$} \Return $\{x_2\}$ \Comment{równanie liniowe}
  \EndIf
  \State $x_1\gets q/A$
  \State \Return $\{x_1,x_2\}$
\Else
  \State $r\gets-B/(2A)$,\quad $s\gets\sqrt{-\Delta}/(2A)$
  \State \Return $\{r+is,\ r-is\}$
\EndIf
\end{algorithmic}
\end{algorithm}

\emph{Poprawność.} Gdy $\Delta\ge0$: $q=-\frac{B+\sqrt\Delta}{2}$, czyli
$q/A=\frac{-B-\sqrt\Delta}{2A}$ to jeden ze szkolnych pierwiastków. Sprawdźmy,
że $C/q$ to drugi --- wystarczą wzory Viète'a. Iloczyn:
$\frac qA\cdot\frac Cq=\frac CA$. Suma: ponieważ $\Delta=B^2-4AC$,
\begin{align*}
q^2&=\frac{B^2+2B\sqrt\Delta+\Delta}{4}=\frac{2B^2+2B\sqrt\Delta-4AC}{4},\\
q^2+AC&=\frac{B(B+\sqrt\Delta)}{2}=-Bq,\\
\frac qA+\frac Cq&=\frac{q^2+AC}{Aq}=-\frac BA.
\end{align*}
Gdy $A=0$, równanie to $Bx+C=0$ i~$x=-C/B$; wtedy $\Delta=B^2$, $q=-B$
i~faktycznie $x_2=C/q=-C/B$. Dzielenie przez $q$ jest bezpieczne, bo
$\abs q\ge B/2>5$. Gdy $\Delta<0$ (możliwe, np.\ $A=C=20$, $B=11$),
musi być $A>0$ i~pierwiastki są zespolone sprzężone.

\emph{Brak catastrophic cancellation.} Na podstawie lematu~\ref{lem:cancel}:
\begin{itemize}
\item $B+\sqrt\Delta$ to suma dwóch liczb nieujemnych --- błąd względny się
nie wzmacnia;
\item $q/A$, $C/q$, $-B/(2A)$, $\sqrt{-\Delta}/(2A)$ to mnożenia, dzielenia
i~pierwiastki --- każde dokłada tylko $\le\eps$ do błędu względnego
(tw.~\ref{tw:model}), błędy względne się \emph{dodają}, nie mnożą przez
duży czynnik;
\item w~przypadku zespolonym część rzeczywistą i~urojoną liczymy osobno,
nic nie odejmujemy.
\end{itemize}
Jedyne odejmowanie to $\Delta=B^2-4AC$. Gdy $B^2\approx4AC$ (pierwiastek
prawie podwójny), tu również tracimy cyfry --- ale to nie wina algorytmu:
wtedy \emph{samo zadanie} jest źle uwarunkowane (zmiana współczynnika
o~$10^{-16}$ przesuwa podwójny pierwiastek o~$\approx10^{-8}$; żaden algorytm
nie da lepszej odpowiedzi z~zaokrąglonych danych). Jeśli zależy nam na
dokładnym $\Delta$, liczymy go w~podwójnej precyzji albo sztuczką Kahana
z~\texttt{fma}: $p=\fl(B^2)$, $d_p=\mathrm{fma}(B,B,-p)$ (dokładny błąd
zaokrąglenia $p$), analogicznie $q'=\fl(4AC)$ i~$d_q$, a~potem
$\Delta=(p-q')+(d_p-d_q)$.

\begin{uwaga}
Założenie $B>10$ gwarantuje, że $B\neq0$ i~znamy jego znak, więc wiemy,
którego znaku użyć (ogólnie: $q=-\tfrac12\bigl(B+\operatorname{sgn}(B)\sqrt\Delta\bigr)$),
oraz że $q$ jest daleko od zera.
\end{uwaga}

\begin{odpowiedz}
Gdy $\Delta=B^2-4AC\ge0$:
\[
q=-\frac{B+\sqrt\Delta}{2},\qquad x_1=\frac qA,\qquad x_2=\frac Cq
\]
(dla $A=0$ tylko $x_2=-C/B$). Gdy $\Delta<0$: $x_{1,2}=\frac{-B}{2A}\pm
i\,\frac{\sqrt{-\Delta}}{2A}$. Nigdzie nie odejmujemy bliskich liczb.
\end{odpowiedz}

% =====================================================================
\section{Zadanie 4. Norma macierzy \texorpdfstring{$\left[\begin{smallmatrix}2&1\\1&2\end{smallmatrix}\right]$}{[[2,1],[1,2]]}}
\begin{tresc}
Na ćwiczeniach zdefiniowaliśmy normę (długość) wektora w~$\R^n$.
Zdefiniujmy \emph{normę macierzy} $A\in\R^{n\times n}$ następująco:
\[
\norm A=\max_{x:\norm x=1}\frac{\norm{A\cdot x}}{\norm x}.
\]
(\ldots czyli „jak bardzo macierz $A$ jest w~stanie wydłużyć wektor”.)
Posługując się tą definicją znaleźć normę macierzy
\[
\begin{bmatrix}2&1\\1&2\end{bmatrix}.
\]
Preferowane są środki dostępne w~tym momencie --- proszę unikać wartości
własnych i~wartości szczególnych.
\end{tresc}

\begin{teoria}[Norma wektora i~norma macierzy]
\begin{definicja}[Norma euklidesowa]
$\norm x=\sqrt{x_1^2+\dots+x_n^2}$ --- zwykła długość wektora.
\end{definicja}

\begin{fakt}
Maksimum w~definicji $\norm A$ istnieje: sfera $\{x:\norm x=1\}$ jest
domknięta i~ograniczona, a~$x\mapsto\norm{Ax}$ ciągła. Ponadto
$\norm{A(tx)}/\norm{tx}=\norm{Ax}/\norm x$, więc to samo wyjdzie, biorąc
maksimum po \emph{wszystkich} $x\ne0$. W~szczególności
$\norm{Ax}\le\norm A\norm x$ dla każdego $x$.
\end{fakt}

Żeby udowodnić $\norm A=M$, trzeba pokazać dwie rzeczy:
\begin{enumerate}
\item \emph{ograniczenie}: $\norm{Ax}\le M$ dla każdego $\norm x=1$;
\item \emph{osiągalność}: istnieje $x$ z~$\norm x=1$ i~$\norm{Ax}=M$.
\end{enumerate}
\end{teoria}

\begin{intuicja}
Macierz jest symetryczna i~„miesza” współrzędne: $A(x_1,x_2)=(2x_1+x_2,\
x_1+2x_2)$. Wektor $(1,1)$ przechodzi na $(3,3)$ --- wydłuża się $3$ razy.
Wektor $(1,-1)$ przechodzi na $(1,-1)$ --- nie zmienia się. Wszystkie inne
kierunki to „mieszanka” tych dwóch, więc wydłużenie jest gdzieś
pomiędzy $1$ a~$3$. Odpowiedź: $3$.
\end{intuicja}

\formalnie
Niech $\norm x=1$, czyli $x_1^2+x_2^2=1$. Wtedy
\begin{align*}
\norm{Ax}^2&=(2x_1+x_2)^2+(x_1+2x_2)^2\\
&=5x_1^2+8x_1x_2+5x_2^2\\
&=5(x_1^2+x_2^2)+8x_1x_2=5+8x_1x_2.
\end{align*}
\emph{Ograniczenie.} Z~$(x_1-x_2)^2\ge0$ mamy $2x_1x_2\le x_1^2+x_2^2=1$, więc
\[
\norm{Ax}^2=5+4\cdot2x_1x_2\le5+4=9,\qquad\norm{Ax}\le3.
\]
\emph{Osiągalność.} Równość jest dla $x_1=x_2$, czyli
$x=\bigl(\tfrac1{\sqrt2},\tfrac1{\sqrt2}\bigr)$:
$Ax=\bigl(\tfrac3{\sqrt2},\tfrac3{\sqrt2}\bigr)$, $\norm{Ax}=3$.

Zatem $\norm A=3$. (Z~parametryzacją $x=(\cos t,\sin t)$ ten sam rachunek
daje $\norm{Ax}^2=5+4\sin2t$, co wprost pokazuje, że wydłużenie zmienia się
od $1$ do $3$.)

\begin{uwaga}
Dla porównania (spoza „dozwolonych środków”): $A$ jest symetryczna z~wartościami
własnymi $3$ i~$1$ (wektory $(1,1)$ i~$(1,-1)$ z~intuicji), a~dla macierzy
symetrycznej norma euklidesowa to największy moduł wartości własnej
--- zgadza się. Dla norm $\norm{\cdot}_1$ i~$\norm{\cdot}_\infty$ normy macierzy to
odpowiednio maksymalna suma modułów w~kolumnie i~w~wierszu --- tu też $3$.
\end{uwaga}

\begin{odpowiedz}
$\norm{\left[\begin{smallmatrix}2&1\\1&2\end{smallmatrix}\right]}=3$,
osiągana dla $x=\frac1{\sqrt2}(1,1)$.
\end{odpowiedz}

\end{document}
